% French PDF pp. 1--6; final sentence completed on p. 7.
\exposetitle{XVIII}{Weil's theorem on the construction of a group from a rational law}
\label{XVIII}
\begin{center}by Michael Artin\end{center}
\renewcommand{\thefootnote}{(\arabic{footnote})}
\footnotetext[0]{Version of 8 May 2009.}

\sgasection{0}{Introduction}
This exposé is devoted to the well-known theorem of Weil [1]
which gives the construction of a group variety from a birational
law. It seems that the generalization of this result to the case
where the base is the spectrum of a discrete valuation ring was
already known to several people; one may, for example, see [2].
Here we shall prove the theorem for a flat group of finite
presentation over an arbitrary base prescheme \(S\).
Since one must make the hypotheses with a little more care,
the statement is not in the form given by Weil; but when \(S\)
is the spectrum of a field, one sees that the form given here
is essentially equivalent to Weil's.
We use the following suggestion of Grothendieck:
given a \og group germ\fg{} \(X\) over a prescheme \(S\),
one must construct the group as the quotient of \(X\times_S X\)
by a suitable equivalence relation.

\sgasection{1}{\og Recollections\fg{} on rational maps}
Let \(X/S\) be a relative prescheme and \(U\subseteq X\)
an open subset. One says that \(U\) is \emph{schematically
dense in \(X\) relative to \(S\)}, or that
\(U\hookrightarrow X\) is \emph{relatively schematically dense},
if for every base change \(S'\to S\) the open subset
\(U\times_S S'\) is schematically dense in \(X\times_S S'\).
For the definition and properties of this notion one refers
to Exp. IX, \(\S\) 4.
{\renewcommand{\thefootnote}{(*)}\addtocounter{footnote}{-1}%
\footnote{Cf. also EGA IV$_3$, 11.9 and 11.10
(in particular 11.10.8), where one says
\og universally schematically dense relative to \(S\)\fg.}}

\begin{proposition}{1.1}\label{XVIII.1.1}
(i) A finite intersection, as well as the union of a nonempty
family of open subsets schematically dense relative to \(S\),
is schematically dense relative to \(S\).

(ii) If \(U\subseteq X\) is schematically dense relative to
\(S\) and \(S\to T\) is a morphism, then
\(U_T\subseteq X_T\) is schematically dense relative to \(T\).

(iii) Let \(U\subseteq V\subseteq X\) be open immersions.
For \(U\) to be schematically dense in \(X\) relative to
\(S\), it is necessary and sufficient that it be so in
\(V\) and that \(V\) be so in \(X\).

(iv) If \(U\subseteq X\) and \(V\subseteq Y\) are relatively
schematically dense, then \(U\times_S V\) is schematically
dense in \(X\times_S Y\) relative to \(S\).
\end{proposition}
In Exp. IX one finds the following criterion:
\begin{proposition}{1.2}\label{XVIII.1.2}
Let \(X/S\) be locally of finite presentation and flat,
and let \(U\) be an open subset of \(X\).
If, for each \(s\in S\), the fiber \(U_s\) is schematically
dense in \(X_s\), then \(U\) is schematically dense in \(X\)
relative to \(S\).
\end{proposition}
In particular:
\begin{corollary}{1.3}\label{XVIII.1.3}
If \(X/S\) is locally of finite presentation and flat,
and each fiber \(X_s\) is \og without embedded components\fg,
then \(U\) is schematically dense in \(X\) relative to \(S\)
if and only if, for each \(s\in S\), \(U_s\) is dense in
\(X_s\) in the topological sense.
\end{corollary}
One easily deduces from the definition the
\begin{proposition}{1.4}\label{XVIII.1.4}
Let \(U\subseteq X\) be schematically dense relative to
\(S\), and let \(f,g:X\to Y\) be two morphisms of functors
over \(S\), where \(Y\) satisfies one of the following conditions:

(i) \(Y\) is a prescheme over \(S\), and each fiber \(Y_s\)
is separated.

(ii) \(Y\) is a presheaf\nde{Specify for which topology:
a priori (fpqc).} over \(S\) such that the diagonal morphism
\(Y\to Y\times_S Y\) is represented by a closed immersion.
\nde{Let us recall here the definition of a closed immersion.
A morphism of \(S\)-functors \(F\to G\) is relatively
representable if, for every \(S\)-morphism \(T\to G\),
where \(T\) is an \(S\)-prescheme, the \(T\)-functor
\(F_T=F\times_S T\) is represented by a \(T\)-prescheme.
A morphism of \(S\)-functors \(F\to G\) is an open
(resp. closed) immersion if it is relatively representable
and, for every \(S\)-morphism \(T\to G\), where \(T\)
is an \(S\)-prescheme, the \(T\)-morphism \(F_T\to T\)
is an open (resp. closed) immersion.}

Then, if \(f=g\) on \(U\), one has \(f=g\).
\end{proposition}
\begin{proof}
In case (i), the standard argument after IX 4.1 shows
that \(f,g\) are equal on each fiber; hence the
subprescheme \(\Ker(f,g)\) of \(X\) is set-theoretically
equal to \(X\), hence is closed, and since it contains
\(U\), it is equal to \(X\), i.e. \(f=g\).
Case (ii) is proved as in \loccit.
\nde{Give details of this point\ldots}
\end{proof}
\begin{definition}{1.5}\label{XVIII.1.5}
{\renewcommand{\thefootnote}{(*)}\addtocounter{footnote}{-1}%
\footnote{Cf. EGA IV$_4$, 20.5, where one says
\og pseudo-morphism from \(X\) to \(Y\) relative to \(S\)\fg,
so as not to conflict with EGA I, 7.12.}}
Let \(X\) be a prescheme over \(S\), and \(Y\) a presheaf
over \(S\). One calls a \emph{rational map}
\(f:X\to Y\) over \(S\) an equivalence class of morphisms
\(f:U\to Y\) over \(S\), where \(U\) is an open subset
of \(X\) schematically dense relative to \(S\), and
where one sets \((f:U\to Y)\sim(f':U'\to Y)\) if and
only if there exists an open subset \(U''\subseteq U'\cap U\),
schematically dense relative to \(S\), such that
\(f=f'\) on \(U''\).
\end{definition}
This definition is made in such a way that one can
define in an evident manner the rational map
\(f\times_S S'\) for every base extension \(S'\to S\).

If \(U\) is an open subset of \(X\), one says that
the rational map \(f\) is \emph{defined on \(U\)}
if there exists a morphism representing \(f\) whose
set of definition contains \(U\).
\begin{definition}{1.5.1}\label{XVIII.1.5.1}
\nde{The number 1.5.1 has been added.}
If \(Y\) satisfies one of conditions (i), (ii) of
1.4, it is clear that there exists a largest open
subset \(U\) of \(X\) such that \(f\) is defined on
\(U\), and this \(U\) is schematically dense relative
to \(S\). One calls it the \emph{domain of definition
of \(f\) over \(S\)}, and denotes it by \(\operatorname{Dom}(f)\).
% typo: duplicated corrupt text "sisisch. dense rel.schématiquement"
% is rendered once as "schematically".
\end{definition}
This notion does not commute with base changes, but one has
\begin{proposition}{1.6}\label{XVIII.1.6}
Let \(X\) be an \(S\)-prescheme and \(Y\) an
\(S\)-functor satisfying one of hypotheses (i), (ii)
of 1.4. Let \(f:X\to Y\) be a rational map over
\(S\), \(S'\to S\) a flat morphism locally of finite
presentation, and \(f'=f\times_S S'\). Then
\[
\operatorname{Dom}(f')=\operatorname{Dom}(f)\times_S S'.
\]
% typo: "Soiten" and "Soenit" -> "Soit".
\end{proposition}
\begin{proof}
Set \(U=\operatorname{Dom}(f)\).
It is clear that \(V'=\operatorname{Dom}(f')\)
contains \(U\times_S S'\).
Let \(V\) be the image of \(V'\), which is an
open subset of \(X\), since \(X'\to X\) is open.
One must prove that \(V=U\), that is, one must
find a morphism \(V\to Y\) representing \(f\).
Set \(S''=S'\times_S S'\),
\[
X''=X\times_S S'',\qquad U''=U\times_S S'',\qquad
V''=V'\times_V V'.
\]
Then \(U''\) is schematically dense in \(X''\)
relative to \(S''\), hence \(U''\) is schematically
dense in \(V''\) relative to \(S\), since
\(U''\subseteq V''\subseteq X''\).
The restriction of \(f':V'\to Y'\) to \(U'\)
is deduced from \(f:U\to Y\) by base change.
The two morphisms \(V''\to Y\) deduced from
\(g\) by base change are equal on \(U''\), hence
equal. Now, since \(V'\to V\) is flat and
locally of finite presentation, it is fppf-covering
(Exp. IV 6.3), and one finds the morphism \(V\to Y\)
by descent.
% typo?: g is not defined; relative to S is kept as printed above.
\end{proof}
We shall frequently use the following triviality:
\begin{proposition}{1.7}\label{XVIII.1.7}
Let \(X/S\) be faithfully flat, locally of finite
presentation, and let \(U\subseteq X\) be a relatively
schematically dense open subset.
Then there exists a base extension \(S'\to S\)
which is fppf-covering, and a section \(x\in X(S')\)
which is contained in \(U(S')\).
\end{proposition}
\begin{proof}
Indeed, \(U/S\) is fppf-covering.
Set \(S'=U\), and take as section the \(S\)-graph
of the inclusion of \(U\) in \(X\).
\end{proof}

\sgasection{2}{Local determination of a morphism of groups}
Let \(G,H\) be groups, and let \(U\subseteq G\)
be a subset such that \(U\cdot U=G\).
Then, if \(f,g\) are homomorphisms from \(G\)
to \(H\) such that \(f=g\) on \(U\), one has \(f=g\).
Similarly:
\begin{proposition}{2.1}\label{XVIII.2.1}
Let \(S\) be a site (cf. Exp. IV\nde{That is,
a category endowed with a topology, cf. IV, \(\S\) 4.2.}),
\(G\) a sheaf of groups on \(S\), and
\(U\subseteq G\) a subsheaf of sets such that
the morphism \(U\times U\to G\) induced by
multiplication is an epimorphism.
Then, if \(f,g\) are homomorphisms from \(G\)
to a sheaf of groups \(H\), equal on \(U\),
one has \(f=g\).
\end{proposition}
\begin{corollary}{2.2}\label{XVIII.2.2}
Let \(G\) be an \(S\)-group prescheme locally
of finite presentation and flat, let
\(U\subseteq G\) be a relatively schematically
dense open subset, and \(H\) a sheaf of groups
on \(S\) for the fppf topology.
Then every homomorphism \(f:G\to H\)
is determined by its restriction to \(U\).
\end{corollary}
\begin{proof}
Indeed, since \(G\) is flat over \(S\),
the composition law in \(G\) is a flat
morphism (VIB.9.2.xi), and it follows that
\(U\times_S U\to G\) is faithfully flat
and locally of finite presentation,
hence fppf-covering, hence an epimorphism.
\end{proof}
\begin{proposition}{2.3}\label{XVIII.2.3}
Let \(G\) be a group prescheme locally
of finite presentation and flat over \(S\),
\(U\subseteq G\) a relatively schematically
dense open subset, and \(H\) a sheaf of
groups for the fppf topology.
Denote by \(m_G\) the multiplication of
\(G\), and by \(m_H\) that of \(H\).

Let \(\overline f:U\to H\) be an \(S\)-morphism,
and suppose that there exists a relatively
schematically dense open subset \(V\) of
\(m_G^{-1}(U)\cap(U\times_S U)\) such that the diagram
\[
\xymatrix{
V\ar[r]^{(\overline f\times\overline f)|_V}\ar[d]_{m_G}
 & H\times_S H\ar[d]^{m_H}\\
U\ar[r]_{\overline f} & H }
\]
is commutative. Then there exists a morphism
of \(S\)-groups \(f:G\to H\) (necessarily
unique) extending \(\overline f\), in each
of the following situations:

(i) \(H\) is representable.

(ii) The diagonal morphism \(H\to H\times_S H\)
is represented by a closed immersion.
\nde{See the note in 1.4.}

(iii) For every section \(a\in U(S)\),
the open subset
\(\mathrm{pr}_2((a\times_S U)\cap V)\)
is relatively schematically dense in \(U\),
{\renewcommand{\thefootnote}{(*)}\addtocounter{footnote}{-1}%
\footnote{Or in \(G\), which amounts to the same
by virtue of 1.1 (iv).}}
and this statement remains true after
every base change \(S'\to S\).
\end{proposition}
\begin{proof}
First note that \(V\) is relatively schematically
dense in \(G\times_S G\).
Indeed, \(V\) is relatively schematically dense in
\[
m_G^{-1}(U)\cap(U\times_S G)\cap(G\times_S U),
\]
and the three factors are deduced from
\(U\subseteq G\) by an evident base change
\(G\times_S G\to G\), which implies the
assertion by 1.1.

To construct a morphism \(f:G\to H\),
it suffices, since \(U\times_S U\to G\)
is fppf-covering, to find a morphism
from \(U\times_S U\) to \(H\) such that
the two induced morphisms on
\((U\times_S U)\times_G(U\times_S U)\)
are the same.
Take \(m_H\circ(\overline f\times\overline f):
U\times_S U\to H\).
One must verify that whenever one has
sections \(a,b,c,d\in U(S')\), with
\(S'\to S\) arbitrary, such that \(ab=cd\),
one also has
\(\overline f(a)\overline f(b)=\overline f(c)\overline f(d)\).
By hypothesis, this is true if
\((a,d)\) and \((c,d)\) are contained in
\(V(S')\), since in this case one has
\(\overline f(a)\overline f(b)=\overline f(ab)\)
and \(\overline f(c)\overline f(d)=\overline f(cd)\).
Thus it is true whenever
\((a,b,c,d)\in(V\times_G V)(S')\).
% typo?: (a,d), rather than (a,b), retained as printed.

I say that \(V\times_G V\) is an open
subset of \((U\times_S U)\times_G(U\times_S U)\)
schematically dense relative to \(S\).
This will complete the proof in cases
(i) and (ii), by 1.4, the facts that the
morphism \(f\) thus constructed extends
\(\overline f\), and that \(f\) is a
homomorphism, being evident, also by 1.4.
Now write
\[
V\times_G V=
\bigl(V\times_G(G\times_S G)\bigr)
\cap\bigl((G\times_S G)\times_G V\bigr).
\]
By symmetry and 1.1, it suffices to
verify that \(V\times_G(G\times_S G)\)
is relatively schematically dense in
\((G\times_S G)\times_G(G\times_S G)\).
But this last prescheme is \(S\)-isomorphic
to \(G\times_S G\times_S G\), the
morphism being given by
\((a,b,c,d)\mapsto(a,b,c)\).
Thus what one must prove is that
\(V\times_S G\) is schematically dense
in \(G\times_S G\times_S G\) relative
to \(S\), which is a consequence of
the fact that \(V\) is relatively
schematically dense in \(G\times_S G\).

It remains to treat case (iii).
To prove
\(\overline f(a)\overline f(b)=\overline f(c)\overline f(d)\),
one may make an fppf-covering base
extension. Suppose that one has a
section \(x\in G(S')\), with
\(S'\to S\) fppf-covering, such that
\((b,x),(a,bx),(d,x),(c,dx)\)
are all in \(V\).
Then one will have
\[
\begin{aligned}
\overline f(a)\overline f(b)\overline f(x)
 &=\overline f(a)\overline f(bx)
=\overline f(abx)=\overline f(cdx)\\
 &=\overline f(c)\overline f(dx)
=\overline f(c)\overline f(d)\overline f(x),
\end{aligned}
\]
whence the desired equality.

\nde{The beginning of the following sentence has been added.}
To find such an \(x\), set, for every \(z\in U(S')\),
\[
V_z=\mathrm{pr}_2((z\times_{S'}U_{S'})\cap V_{S'}).
\]
Then the hypotheses on \(x\) mean that \(x\in W\), where
\[
W=V_b(S')\cap V_d(S')\cap b^{-1}V_a(S')\cap d^{-1}V_c(S').
\]
By (iii), \(W\) is relatively schematically
dense in \(G\).
Thus the existence of a section \(x\)
after an fppf-covering extension follows
from Proposition 1.7.

Let us verify that the morphism thus
constructed is multiplicative:
let \(a,b\in G(S')\), with \(S'\to S\)
arbitrary, and write \(S'=S\).
Choose an fppf-covering base extension
\(S'\to S\) and a section \(x\in G(S')\)
such that \(x\in U(S')\) and
\(ax^{-1}\in U(S')\), that is,
\(x\in(a^{-1}U)^{-1}(S')\).
Choose moreover another fppf-covering
extension \(S''\to S'\) and a section
\(y\in G(S'')\) such that
\[
y,\;by^{-1},\;aby^{-1}\quad\text{belong to }U(S''),
\]
and
\[
by^{-1}\in V_x,\qquad xby^{-1}\in V_{ax}^{-1}.
\]
Then, by definition of \(f\), one has over \(S''\)
\[
\begin{gathered}
f(a)=\overline f(ax^{-1})\overline f(x),\qquad
f(b)=\overline f(by^{-1})\overline f(y),\\
f(ab)=\overline f(aby^{-1})\overline f(y).
\end{gathered}
\]
Moreover,
\[
\begin{aligned}
f(a)f(b)
 &=\overline f(ax^{-1})\overline f(x)
\overline f(by^{-1})\overline f(y)\\
 &=\overline f(ax^{-1})\overline f(xby^{-1})\overline f(y)\\
 &=\overline f(aby^{-1})\overline f(y)=f(ab),
\end{aligned}
\]
whence multiplicativity.
% typo?: V_ax^{-1} and the adjacent V_x condition are retained as printed.

The fact that \(f\) extends \(\overline f\)
is now easy: let \(a\in U(S')\), write
\(S'=S\), and choose \(S'\to S\)
fppf-covering and sections
\(x,y\in G(S')\) such that
\((x,y),(ax,y),(a,xy)\) are in \(V(S')\).
Then
\[
\begin{gathered}
f(xy)=\overline f(x)\overline f(y)=\overline f(xy),\\
f(axy)=\overline f(ax)\overline f(y)=\overline f(axy).
\end{gathered}
\]
Thus
\[
\begin{aligned}
f(a)&=f(axy)f((xy)^{-1})
=f(axy)f(xy)^{-1}\\
&=\overline f(axy)\overline f(xy)^{-1}=\overline f(a).
\end{aligned}
\]
\end{proof}
\begin{remark}{2.3.1}\label{XVIII.2.3.1}
In many cases hypothesis (iii) is true,
for it will even be true if one
replaces \(U\) by a smaller open subset
\(U'\), still relatively schematically
dense in \(G\). For example, one has:
\end{remark}
\begin{proposition}{2.4}\label{XVIII.2.4}
With the situation as in 2.3, suppose
that each geometric fiber of \(G/S\)
is irreducible.
Let \(U'=\mathrm{pr}_1 V\), and
\[
V'=V\cap m_G^{-1}(U')\cap(U'\times_S U').
\]
Then \(U'\subseteq U\) is a relatively
schematically dense open subset of \(G\),
and the objects \(U'\), \(\overline f|_{U'}\),
and \(V'\) satisfy hypothesis (iii).
\end{proposition}
\begin{proof}
\(U'\) is open, since \(G\times_S G\to G\)
is flat and locally of finite presentation.
All the other verifications are trivial,
except for hypothesis (iii).
Let \(a\in U(S')\), and write \(S'=S\).
To verify that
\(\mathrm{pr}_2((a\times U')\cap V')\)
is relatively schematically dense in
\(U'\), it suffices to do so fiber
by fiber by Corollary 1.3, that is,
it suffices to treat the case where
\(S\) is the spectrum of a field,
and in this case it suffices to prove
that \(U'\) is nonempty, since \(G\)
is irreducible and \og without embedded
components\fg{} (cf. VIA, 1.1.1).
Now
\[
\begin{aligned}
\mathrm{pr}_2((a\times U')\cap V')
&=\mathrm{pr}_2((a\times U')\cap V)\\
&\quad\cap\mathrm{pr}_2((a\times U')\cap m_G^{-1}(U')),
\end{aligned}
\]
and the second term of the right-hand
member is dense in \(G\).
% Proof continues in en-02.tex.
